% =============================================================
% CHAPTER 4 — METHODOLOGY AND EXPERIMENTAL FRAMEWORK
% =============================================================
\chapter{Methodology and Experimental Framework}
\label{chap:methodology}

This chapter presents the methodology and experimental framework designed to benchmark and evaluate Multimodal Large Language Models (MLLMs) for zero-shot and few-shot face-swap detection. Rather than fine-tuning models on task-specific datasets, this work evaluates the intrinsic visual forensic capabilities of off-the-shelf instruction-tuned MLLMs under controlled prompting regimes.

The remainder of this chapter is organised as follows:
\Cref{sec:formulation} defines the formal task formulation and deterministic response-parsing protocol;
\Cref{sec:datasets} details the benchmark datasets, manipulation generators, and the balanced development set;
\Cref{sec:models_infra} presents the evaluated model suite, the inference pipeline, and GPU memory isolation architecture;
\Cref{sec:staged_prompting} details the staged zero-shot prompt engineering funnel;
\Cref{sec:few_shot_method} describes the few-shot in-context learning experimental design; and
\Cref{sec:eval_protocol} outlines the evaluation protocol, performance metrics, and reproducibility environment.

% -------------------------------------------------------------
\section{Problem Formulation and Response Parsing}
\label{sec:formulation}

\subsection{Mathematical Task Formulation}

We formulate single-image face-swap detection as a binary classification problem. Let $\mathcal{X}$ denote the space of RGB facial images. An input image $x \in \mathcal{X}$ is associated with a ground-truth label $y \in \mathcal{Y} = \{\text{real}, \text{fake}\}$, where:
\begin{itemize}
    \item $\text{real}$ ($y = 0$): corresponds to an authentic, unaltered facial photograph of a person;
    \item $\text{fake}$ ($y = 1$): corresponds to a digitally manipulated image in which the facial identity of the original subject has been replaced by a target identity via an automated face-swapping algorithm.
\end{itemize}

Let $m$ denote an off-the-shelf instruction-tuned MLLM. Given an input facial query image $x$, a prompt $p$, and a hyperparameter configuration $\theta_{\text{gen}}$ (such as sampling temperature and maximum output sequence length), the model autoregressively generates an output text sequence $t$:
\begin{equation}
    t = m(p, x; \theta_{\text{gen}})
\end{equation}

\subsection{Prompt Parameterization}

In typical applications, prompts are formulated as monolithic text strings. However, modifying multiple prompt aspects simultaneously makes it impossible to isolate which specific factor drives changes in model performance. To enable controlled and systematic evaluation, we formulate a prompt $p$ as a modular tuple of five distinct components:
\begin{equation}
    p = \langle \mathcal{V}, \mathcal{S}, \mathcal{R}, \mathcal{C}, \mathcal{D} \rangle
\end{equation}
where:
\begin{itemize}
    \item $\mathcal{V} = (w_{\text{real}}, w_{\text{fake}})$ specifies the \textbf{target class vocabulary} words the model is instructed to output for real and fake faces (e.g., \texttt{("authentic", "faceswapped")});
    \item $\mathcal{S}$ defines the \textbf{extra context} provided in the system or user prompt;
    \item $\mathcal{R}$ specifies the \textbf{persona} in the systemp prompt (e.g., \emph{Forensic Video Analyst} );
    \item $\mathcal{C}$ governs the \textbf{reasoning mode} determining whether and how the model should produce a reasoning for it's choice (e.g. Chain-of-Thought);
    \item $\mathcal{D} = \{(x_1^{(d)}, y_1^{(d)}), \dots, (x_n^{(d)}, y_n^{(d)})\}$ denotes the \textbf{in-context demonstration set} of $n$ labelled visual exemplars ($n = 0$ in zero-shot configurations).
\end{itemize}

This decomposition directly enables the staged experimental funnel described in \Cref{sec:staged_prompting} and the few-shot protocol in \Cref{sec:few_shot_method}.

\subsection{Deterministic Parsing and Failure Semantics}

Because MLLMs emit natural language, benchmarking their predictions requires a deterministic mapping from the generated token sequence $t$ to the binary label space $\mathcal{Y} = \{\text{real}, \text{fake}\}$. We define a deterministic parser $\pi(t)$ that maps the generated output to a predicted label according to the declared extraction policy:
\begin{equation}
    \hat{y} = \pi(t) \in \{\text{real}, \text{fake}, \bot\}
\end{equation}
where $\bot$ denotes a distinguished \textbf{parse-failure} symbol emitted when the model fails to produce an output that can be mapped to either \emph{real} or \emph{fake} (e.g., due to model refusals or not following the prompt instructions).

% -------------------------------------------------------------
\section{Benchmark Datasets}
\label{sec:datasets}

To evaluate how well the models generalise across different manipulation algorithms, subjects, and occlusion conditions, we use two video benchmarks for face swapping: \textbf{GOTCHA}~\cite{gotcha} and \textbf{FOWS} (\emph{Face Occlusion With Swapping})~\cite{fows}. In both, subjects perform physical actions that cover part of the face, producing frames that stress face swapping generators. Each frame is accordingly labelled as occluded (\texttt{occ}) or non-occluded (\texttt{no-occ}), depending on whether the hand or object covers the face.

\subsection{Source Datasets}

\paragraph{GOTCHA.}
GOTCHA~\cite{gotcha}, introduced by Mittal et al., is a benchmark for challenge-response detection of real-time deepfakes, in which subjects perform prescribed actions that expose the limitations of live face swapping pipelines. It contains recordings of 47 subjects performing eight challenges grouped into four categories: head movement, face occlusion, facial deformation, and face illumination, with swaps available from three generators. Following Ziglio et al.~\cite{fows}, we use the hand and object occlusion challenges and the swaps produced by DeepFaceLab~\cite{deepfacelab} and FSGAN (v2)~\cite{fsgan}.

The 47 subjects are partitioned into 30 for training and 17 for testing, with no subject appearing in both sets as either source or target identity. All evaluation uses the 17-subject \textbf{test partition}; the 30 training subjects serve only as the few-shot demonstration pool (\Cref{sec:few_shot_method}). Over the hand and object occlusion challenges, the test partition comprises \num{68899} frames (\num{36236} real and \num{32663} fake).

\paragraph{FOWS.}
FOWS (Face Occlusion With Swapping) is a dataset collected by Ziglio et al.~\cite{fows} at Fondazione Bruno Kessler (FBK). It focuses on the face occlusion setting, which prior work identified as the most effective at revealing face swapping artifacts in remote video communication. Each clip shows a subject who, on request, occludes their own face either with their hand (hand occlusion challenge) or with a supplied rectangular object (object occlusion challenge) at predefined positions.

Recording was driven by a guiding video overlaid on the camera feed that indicated where to place the face and the occluder. For each challenge they defined three variants with different occluder trajectories, for six guiding videos of about 40 seconds. Every clip starts with a static segment showing the face without occlusion and continues with a dynamic segment of roughly 15 seconds in which the occluder moves across the face. Seven subjects each recorded the six variants, giving 42 genuine clips.

For each genuine recording, manipulated counterparts were produced with three generators that swap faces from a single target image, SimSwap~\cite{simswap}, GHOST~\cite{ghost}, and FaceDancer~\cite{facedancer}, using royalty-free celebrity photographs as target identities. A fourth generator, DLC, was added after publication and is not described in the paper.
% TODO: add citation and full name for DLC once located.

In total, FOWS provides \num{42000} frames over the seven subjects, evenly split between occluded and non-occluded frames. Five subjects form the training pool used for few-shot demonstrations (\Cref{sec:few_shot_method}); the remaining two form the \textbf{test partition}, a deliberately challenging split of the only female subject and the only male subject with a thick beard. The test partition contains \num{12000} frames (\num{2400} real and \num{9600} fake), a $1:4$ ratio of real to fake frames obtained by pairing each genuine frame with its four swaps.

\Cref{tab:datasets} reports the full composition of both test partitions, broken down by source (genuine frames or individual generator) and by occlusion condition.

\begin{table}[t]
	\centering
	\caption{Composition of the evaluation test partitions by source and occlusion condition (frame counts). Rows give the genuine frames and each face swapping generator; training subjects are excluded and used only as the few-shot demonstration pool.}
	\label{tab:datasets}
	\begin{tabular}{lrrr}
		\toprule
		\textbf{Source} & \textbf{Occluded (\texttt{occ})} & \textbf{Non-occluded (\texttt{no-occ})} & \textbf{Total} \\
		\midrule
		\multicolumn{4}{l}{\textit{GOTCHA} (17 test subjects)} \\
		\midrule
		Genuine        & \num{23000} & \num{13236} & \num{36236} \\
		DeepFaceLab    & \num{9899}  & \num{6338}  & \num{16237} \\
		FSGAN (v2)     & \num{10019} & \num{6407}  & \num{16426} \\
		\cmidrule(lr){1-4}
		Fake subtotal  & \num{19918} & \num{12745} & \num{32663} \\
		\textbf{Total} & \num{42918} & \num{25981} & \num{68899} \\
		\midrule
		\multicolumn{4}{l}{\textit{FOWS} (2 test subjects)} \\
		\midrule
		Genuine        & \num{1200} & \num{1200} & \num{2400} \\
		SimSwap        & \num{1200} & \num{1200} & \num{2400} \\
		GHOST          & \num{1200} & \num{1200} & \num{2400} \\
		FaceDancer     & \num{1200} & \num{1200} & \num{2400} \\
		DLC            & \num{1200} & \num{1200} & \num{2400} \\
		\cmidrule(lr){1-4}
		Fake subtotal  & \num{4800} & \num{4800} & \num{9600} \\
		\textbf{Total} & \num{6000} & \num{6000} & \num{12000} \\
		\bottomrule
	\end{tabular}
\end{table}

\subsection{Development Set for Prompt Selection}
\label{subsec:devset}

The staged prompt search (\Cref{sec:staged_prompting}) and the few-shot configuration search (\Cref{sec:few_shot_method}) each require evaluating a large number of candidate configurations. Running every candidate over the full test partitions, which together exceed \num{80000} frames, is computationally prohibitive. We therefore construct a compact development set for each dataset, use it to select configurations during the search phases.

Each development set is a random subsample of the corresponding test partition, balanced along every axis present in the data. \Cref{tab:devbalance} reports the frame count in each category of every axis.

\begin{table}[t]
	\centering
	\caption{Balance of each development set along its axes. A dash marks a generator the dataset does not use. For FOWS, each occluder challenge is split equally across three movement paths (\num{224} frames each).}
	\label{tab:devbalance}
	\begin{tabular}{llrr}
		\toprule
		\textbf{Axis} & \textbf{Category} & \textbf{GOTCHA} & \textbf{FOWS} \\
		\midrule
		Class label          & Real                        & \num{672} & \num{672} \\
		                     & Fake                        & \num{672} & \num{672} \\
		\midrule
		Occlusion condition  & Covered (\texttt{occ})      & \num{672} & \num{672} \\
		                     & Uncovered (\texttt{no-occ}) & \num{672} & \num{672} \\
		\midrule
		Generator            & DeepFaceLab & \num{336} & --        \\
		(fake frames)        & FSGAN       & \num{336} & --        \\
		                     & SimSwap     & --        & \num{168} \\
		                     & GHOST       & --        & \num{168} \\
		                     & FaceDancer  & --        & \num{168} \\
		                     & DLC         & --        & \num{168} \\
		\midrule
		Occluder challenge   & Hand occlusion   & \num{672} & \num{672} \\
		                     & Object occlusion & \num{672} & \num{672} \\
		\midrule
		Subject              & Number of test subjects & 14       & 2         \\
		                     & Frames per subject      & \num{96} & \num{672} \\
		\bottomrule
	\end{tabular}
\end{table}

Both development sets contain \num{1344} frames (\num{672} real, \num{672} fake). The FOWS set covers both test subjects; the GOTCHA set covers 14 of the 17 test subjects, the remaining three lacking enough frames in some axis to fill the balanced quota.

\guide{TODO: Justify why the development set is balanced along every axis instead of preserving the natural test-partition prevalence }





